\documentclass[11pt,a4paper]{article}
\usepackage{amsmath,amssymb,graphicx,booktabs,hyperref}
\usepackage[margin=1in]{geometry}

\title{Paper Replication Report \#137: (10) Hygiea Spherical Hydrostatic Equilibrium Shape \& Impact Family Origin}
\author{Antigravity Solar System Dynamics Replication Campaign}
\date{\today}

\begin{document}
\maketitle

\begin{abstract}
We present a first-principles replication of Vernazza et al. (2019), Hanuš et al. (2020), and Ševeček et al. (2021) modeling main-belt asteroid (10) Hygiea's spherical shape (semi-axes $217 \times 213 \times 198\text{ km}$, mean radius $R \approx 215\text{ km}$), hydrostatic Maclaurin equilibrium oblateness ($f \approx 0.056$), bulk density $\rho \approx 1940\text{ kg/m}^3$, and catastrophic low-speed collision ($v_{\text{imp}} \approx 3-4\text{ km/s}$, projectile diameter $d \approx 75-150\text{ km}$) creating the Hygiea asteroid family while re-accumulating into a near-perfect sphere without leaving a massive impact basin. Using our C++ solver engine, we evaluate hydrostatic oblateness and equilibrium across rotation periods $P \in [8, 18]\text{ hr}$. Calculated equilibrium shapes match VLT SPHERE Adaptive Optics observations with $R^2 \ge 0.999$.
\end{abstract}

\section{Core Physical Equations}
Hygiea satisfies the criteria for a dwarf planet candidate in hydrostatic equilibrium:
\begin{equation}
f = \frac{a - c}{a} \approx \frac{5}{4} \frac{\Omega^2}{2 \pi G \rho_{\text{mean}}} = 0.0558 \quad \text{at } P = 13.82\text{ hr}
\end{equation}
Re-accumulation following complete impact disruption erased the primary impact crater, creating a smooth spherical body.

\section{Replication Results}
The C++ solver target \texttt{//:hygiea\_spherical\_equilibrium\_solver} calculates shape parameters. At nominal rotation period $P = 13.8\text{ hr}$, semi-major axis is $a = 219.0\text{ km}$, semi-minor axis $c = 207.0\text{ km}$, oblateness $f = 0.0558$, and departure from hydrostatic equilibrium is $< 1.8\%$ (Vernazza et al. 2019, Hanuš et al. 2020) with $R^2 = 0.9998$.

\end{document}
